\subsection{Full-Inference DeepSeek: One Chip, Prefill and Decode}
\label{sec:decode_decomp}\label{sec:batch_sensitivity}
We evaluate DeepSeek-V3 on 32 independent requests, each with 4,096 prompt
and 512 output tokens. All 61 layers execute: prefill produces the first
output, followed by 511 decode forwards. Captured expert IDs and routing
weights define all $61{,}985{,}244$ invocations in $41{,}924{,}510$ execution
groups. Original FP8 weights are converted to BF16 with the captured rounding
policy; higher-precision tensors and complete KV history are preserved.

The recovered chip has 320 $16\times16$ MAC arrays, 5,248 vector lanes,
1,024 paid FP8 receiver lanes, 78.25 MiB SRAM, 64 GiB DRAM and 2 TiB backing.
At a 1 ns analytical clock, the corrected chip provides 1,630 GB/s DRAM,
640 GB/s external and 20,480 GB/s local service. It passes the shared BF16
scenario's area, peak-power, density, pins, bank, depth and NoC predicates:
$321.69$ mm$^2$ and $298.69$ W. The BF16 peak coefficient is 1.2 mW/MAC;
the area scale and interface organization remain explicit concept assumptions,
not calibrated silicon estimates. Retiming preserves the original cover,
addresses and outer dependency order; 384 regenerated pipeline recipes pass
replay, and all other durations are conservatively extended.
The historical TPU-v3 extension retains the reference compute/DRAM resources
with 32 GB/s external and 8,192 GB/s local service, and 2,048 primitive vector
cycles/ns. These interface assumptions are explicit parts of the prediction,
not measurements of a calibrated TPU serving stack or fabricated custom chip.
This extension is not yet a verified reconstruction of FAST's modeled TPU-v3
baseline, which includes ordinary XLA fusion. The ratios below use that
retained reference with the corrected custom chip; they remain conditional
comparisons, not established reproduction of FAST's modeled TPU-v3 baseline.

\begin{table}[t]
\centering\small
\caption{Corrected custom-chip predictions (seconds); ratios use the retained, uncalibrated TPU-v3 extension.}
\label{tab:native_phases}
\begin{tabular}{@{}lrrr@{}}
\toprule
Phase & TPU-v3 extension & Custom chip & Speedup \\
\midrule
Prefill & 13,387.193 & 1,078.914 & $12.408\times$ \\
Decode & 53,083.324 & 2,074.999 & $25.582\times$ \\
Preparation & 131.252 & 42.973 & $3.054\times$ \\
\midrule
Total & 66,601.769 & 3,196.886 & $\mathbf{20.833\times}$ \\
\bottomrule
\end{tabular}
\end{table}

Phase times add; speedup is the ratio of total predicted inference latencies.
Decode contributes $80.45\%$ of time saved, prefill $19.41\%$, and preparation
$0.14\%$. Initial conversion is charged in preparation; subsequent streamed
conversion is charged to the consuming kernels. Optimization time is separate.
The retained execution's state replay is reused with the checked retiming;
the corrected fixed-chip LP separately has a $0.0846842\%$ objective interval
(Appendix~\ref{app:corrected_native_lp}). Neither this interval nor the
retiming establishes global integral optimality.

\subsection{Feature Contributions on the Same Chip}
\label{sec:native_features}
Table~\ref{tab:native_features} retains the archived chip and removes one
mechanism at a time on the original numerical workload. These are matched
execution controls, not globally optimized feature-off designs. Interacting
rows are not an additive attribution of the headline speedup. They have not
been retimed to the corrected chip in Table~\ref{tab:native_phases}.

\begin{table}[t]
\centering\small
\caption{Archived-chip feature controls: predicted time saved (seconds).}
\label{tab:native_features}
\begin{tabular}{@{}lrr@{}}
\toprule
Feature & Prefill & Decode \\
\midrule
Matrix batching & 43.173 & 856.070 \\
SwiGLU primitive fusion & 14.579 & 1.819 \\
Other elementwise register storage & 257.071 & 34.026 \\
Full-K panel reuse & 32.445 & 0 \\
Per-use weighted decode tiles & 0 & 0.866 \\
Optional KV residency & 0 & 15.726 \\
Phase-reused placement (included above) & 0 & 3.698 \\
\bottomrule
\end{tabular}
\end{table}

\paragraph{Batching and storage fusion.}
The no-batch control executes every original member separately and reselects
its tile. Resident weights remain resident; each streamed member pays its own
receive. The control expands to exactly $61{,}985{,}244$ kernels and takes
$2{,}999.864$ seconds, versus $2{,}100.621$ with batching. Thus grouped execution
and shared streamed-weight delivery give $1.612\times$ decode and
$1.064\times$ prefill gains. SwiGLU and elementwise storage choices remove
primitive spills while retaining arithmetic, BF16 rounding, and local/vector
service. Full-K panels reuse weights across output-row tiles, paying SRAM and
serialized panel waves. Holding one tile fixed per parameter family adds $0.866$ seconds to decode;
panel boundaries, buffers and residency stay fixed. This tests weighted decode
tile adaptation, rather than every temporal choice in the master.

\paragraph{KV and time-indexed placement.}
Every arm retains complete cached decoding, growth, host backing and mandatory
writes. Optional DRAM residency selects 246 request/layer/K-or-V prefix chains
and pays 125,952 copies. A fixed whole-task partition fits 193 chains; releasing
prefill temporaries before decode permits 246. Placement's $3.698$-second gain
is included in the $15.726$-second KV gain. These modest marginal gains reflect
limited DRAM capacity and the wide custom external interface. Independent
checks replay changed copies and kernel traffic; Appendix~\ref{app:native_provenance}
records exact inputs, costs, controls and their scope.
